\section{R19: Primitive residual incidence and its limits}
\label{sec:r19-partial}
\noindent\textbf{Research status.} The joint primitive theorem sought in
R19 is not established. We retain the conditional mechanism and its
original negative cases. The following result is a proved partial
comparative static, not a claim about every demand or competition game.

\subsection{Consumer primitives and a common post-adoption component}
Consider baseline utility
\(U_0(q)=\mathbf1^\top q-q^\top M_0q/2\), where
\(M_0=(1-\gamma)I+\gamma\mathbf1\mathbf1^\top\).
Buyers take expected network benefits as given individually. Fully
compatible fulfilled networks give
\(p=\mathbf1-(M_0-v\mathbf1\mathbf1^\top)q\) and
\(CS=q^\top M_0q/2\). At \(\gamma=1\), retain the original
uniform-type tied-buyer allocation and quantity clearing, rather than
claiming strict concavity of differentiated demand.

Two symmetric member suppliers have zero marginal cost, and the outsider
has residual cost \(d\). The separately derived own-markup coefficients
are \(\ell_C=1-v\) for Cournot, and
\[
\ell_B=\frac{(1-\gamma)(1+2\gamma-3v)}{1+\gamma-2v}
\]
for differentiated Bertrand. These follow from the actual quantity and
price strategy equations respectively; an arbitrary \(\ell\) is only an
algebraic conduct coefficient, not an economic primitive.

Put \(b=\gamma-v\ge0\), \(\tau=1-\gamma+\ell>0\), and
\(\kappa=\tau+3b>0\). The fully active solution is
\[
x(d)=\frac{\tau+bd}{\tau\kappa},\qquad
 y(d)=\frac{\tau-(\tau+2b)d}{\tau\kappa},\qquad
0\le d<\bar d=\frac{\tau}{\tau+2b}.
\]
The member national component from its two symmetric member destinations
is
\[
m_C(d)=(1+\gamma+2\ell)x(d)^2+2\gamma x(d)y(d)+y(d)^2/2.
\]
The two profit terms represent the \emph{same home firm's} profits in
two destinations by symmetry. They do not add foreign profits to welfare.

\begin{proposition}[Primitive residual-incidence boundary, R19 T-R]
\label{prop:r19-residual}
Define
\begin{align*}
\rho&=1+2\gamma-4\ell b/\tau,\\
\mathcal A&=\tau^2+4(1-\gamma)\tau b+
 [6(1-\gamma)+4\ell]b^2>0.
\end{align*}
On the stated active branch,
\[
R_{\gamma,v,\ell}(d):=m_C(d)-m_C(0)
=\frac{-\rho d+\mathcal A d^2/(2\tau^2)}{\kappa^2}.
\]
It is strictly convex. If \(\rho\le0\), it is nondecreasing for
\(d\ge0\). If \(\rho>0\), its unconstrained minimum is
\(d_{\min}=\rho\tau^2/\mathcal A\), and it decreases before this
point. Every turning or zero-return point must be truncated by
\(d\le c\), \(d<\bar d\), and the actual investment equilibrium.
For original homogeneous Cournot, \(\rho=4v-1<0\) on \(v<1/4\).
For substitute differentiated Bertrand, \(0\le v\le\gamma<1\)
implies \(\rho>0\).
\end{proposition}
\begin{proof}
The symmetric system
\((M_0-v\mathbf1\mathbf1^\top+\ell I)q
=\mathbf1-(0,0,d)^\top\) gives the displayed quantities. Expanding
\(m_C(d)-m_C(0)\) gives the linear coefficient
\(-\rho/\kappa^2\) and quadratic coefficient
\(\mathcal A/(2\tau^2\kappa^2)\).
Positivity follows term by term. At \(\gamma=1,\ell=1-v\),
this is identically the original equation~\eqref{eq:residual-rent}.
For Bertrand substitution gives
\[
\rho=\frac{2\gamma^2-10\gamma v-3\gamma+12v^2+v-2}
{-3\gamma+5v-2}.
\]
The denominator is negative. Its numerator is convex in \(v\) and is
negative at both endpoints \(v=0,\gamma\): respectively
\((\gamma-2)(2\gamma+1)\) and
\(2(\gamma-1)(2\gamma+1)\). Convexity bounds it above by the
negative chord. Thus \(\rho>0\). Differentiating the quadratic
proves the remaining assertions.
\end{proof}

\subsection{Global price continuation and endogenous-investment limits}
For the fully compatible price game, write
\(M=(1-\gamma)I+b\mathbf1\mathbf1^\top\).
It is positive definite for \(\gamma<1,b\ge0\), giving unique
consumer demand at every price by complementarity. As an active firm's
price rises, rivals only enter; with \(n\) active brands its demand slope
magnitude is
\[
\frac{1-\gamma+(n-1)b}{(1-\gamma)(1-\gamma+nb)}.
\]
This magnitude increases at rival entry. Thus its positive-sales profit
is piecewise strictly concave with downward derivative jumps, and its
interior stationary price is a global best reply. Below-cost or no-sale
prices cannot dominate its strictly positive profit.

Every zero-cost member must enter a putative boundary equilibrium. If
only the outsider were inactive, including at its choke-price kink,
necessary right-derivative conditions imply member markups
\(p_i\ge\ell_Bq_i\), and hence
\(Q\le2/(\tau+2b)\). The outsider's reservation price is then at
least \(\tau/(\tau+2b)=\bar d>d\), admitting profitable entry.
All equilibria must therefore be active, and the linear active FOC system
has a unique solution. This establishes unique \emph{pure-price} Nash
in this full-compatible family. Other portfolio network matrices require
separate checks; full compatibility alone does not certify them.

Holding the pre-adoption member-market outsider profit \(B\) fixed,
its privately financed two-destination support gain is
\(G_o(d)=2[\ell y(d)^2-B]\), with
\[
G_o'(d)=-\frac{4\ell(\tau+2b)y(d)}{\tau\kappa}<0.
\]
When \(\rho>0\) and \(d<d_{\min}\), raising residual cost
therefore \emph{weakens private adoption while strengthening the member's
relative IS preference}, conditional on strict outsider-only adoption
remaining selected. The outsider's strict best-response cutoff is
\[
d_A=\bar d\left[1-\kappa\sqrt{(B+F/2)/\ell}\right]
\]
when positive; it does not replace all member and SW investment checks.

At the R16 Bertrand primitives
\((\gamma,v,c,d,F)=(23/25,33/500,2/125,1/500,1/40)\),
exact investment checks give a strict private interval. Government gaps are
\[
\Delta_0\simeq+0.00277401,\qquad \Delta_1\simeq-0.000907434.
\] These are different primitives from the original
Cournot example. Here \(d_A\simeq0.0105895<c=0.016<d_{\min}
\simeq0.040003\); no observable U-shape or later SU return is claimed.
Original R10's identical-parameter full reversal remains withdrawn.

For differentiated Cournot with \(0\le v<1/4\), the residual first
slope switches at
\[
\gamma_C^*(v)=\frac{2v-1+\sqrt{17+20v-28v^2}}4.
\]
Indeed \(\rho_C=N_C/(\gamma+v-2)\), with
\(N_C=2\gamma^2-2\gamma v+\gamma+4v^2-3v-2\).
This numerator is strictly increasing in \(\gamma\ge v\), negative
at \(\gamma=v\) and positive at \(\gamma=1\). Its unique root is
the displayed cutoff, giving a negative residual slope below it and a
positive one above it. Thus substitution and markups jointly determine
the contrast; it is not a universal conduct dichotomy.

\subsection{What remains unproved and what is already known}
One-way private investment does not imply a government ranking reversal:
with original Cournot primitives \((v,c,d,F)=(0,1/10,1/20,1/50)\),
all-profile strict dominance gives outsider-only SU support and no SW
support, but the initial government gap is \(-47/3200\).
R19 has not derived the initial SU advantage and all investment threshold
orders from a single common primitive condition. Overall grade therefore
remains Level~2; the cumulative Level~3 and~4 criteria are not met.

General welfare incidence and inverse-Jacobian methods are already
established by \citet{r19WeylFabinger2013,r19AdachiFabinger2022}. Public standards
followed by private fixed-cost adaptation are already in Maggi and
Mr\'azov\'a; one-way compatibility and both competition forms are prior, including \citet{r19BuccellaEtAl2026}.
The precise institutional residual increment remains subject to source
comparison, with latest-2026 and Walz--Woeckener full-proof gaps. No
firstness claim is made. The nonlinear Bertrand examples remain numerical
checks, not a global nonlinear-price uniqueness theorem.
